\section{Frozen-model transfer to Mureka500}
\label{sec:mureka-frozen-transfer}

The 500 source-labelled Mureka v9 recordings were evaluated using the
previously stored Human/Suno development models, with no fitting, calibration,
threshold adjustment or source-dependent preprocessing. The same 500 native
center-60-second intervals were used for every candidate. All 127 nonempty
S/D/R/P/F/H/M subsets, five training-group caps and five original fold models
produced 3,175 model instances and 1,587,500 prediction rows. These rows are
dependent applications to 500 songs, not independent audio trials.

The independent numerical auditor reconstructed every score using the saved
training medians, missing indicators, scales and coefficients. All decisions
and all 3,175 sensitivity summaries agreed. Maximum score discrepancy was
$1.9984\times10^{-15}$; the nearest independent score was
$2.5442\times10^{-7}$ from the fixed threshold of 0.5. No missing-feature row
was excluded: 15 recordings lack all six P descriptors and one lacks all six
M descriptors. Their original training-derived missing-value treatment remains
unchanged. Raw linear scores are not calibrated probabilities.

\subsection{All-cap comparison}
Here ``all'' means all available training groups within each original fold,
not training a new model on all 1,604 development recordings. The five stored
training folds contain 1,250--1,328 rows. Each model is applied to all 500 new
recordings. Table~\ref{tab:mureka-main} reports mean sensitivity across the
five stored models. Their ranges are descriptive, not confidence intervals.

\begin{table}[htbp]
\centering\small
\begin{tabular}{lrrr}
\hline
Feature set & Sensitivity (\%) & FNR (\%) & Change (pp)\\
\hline
S/D/R/P & 76.76 & 23.24 & ---\\
S/D/R/P + F & 63.76 & 36.24 & $-13.00$\\
S/D/R/P + H & 80.08 & 19.92 & $+3.32$\\
S/D/R/P + M & 83.28 & 16.72 & $+6.52$\\
S/D/R/P/F/H/M & 79.36 & 20.64 & $+2.60$\\
\hline
\end{tabular}
\caption{Frozen transfer to Mureka500 at the original all-cap operating point.
Each percentage is the mean over five dependent model applications, each with
denominator 500. Change is relative to S/D/R/P.}
\label{tab:mureka-main}
\end{table}

The baseline fold sensitivities range from 66.20 to 86.40\%; the corresponding
ranges are 48.00--74.20\% with F, 72.20--86.60\% with H, 73.40--90.20\%
with M, and 71.20--87.40\% with all seven. F decreases sensitivity in every
all-cap fold comparison, while M increases it in every such comparison.
This is a descriptive paired result, not a significance test or selection of
a universally superior candidate.

\subsection{Training-size and single-family observations}
All-cap single-family mean sensitivities are S 89.24\%, D 48.32\%, R 91.00\%,
P 67.12\%, F 64.20\%, H 80.56\% and M 82.12\%. High AI sensitivity alone
does not establish Human specificity. The fixed test denominator remains 500
under every training cap; a cap is an original per-source training-group budget,
not that number of test songs or necessarily that number of training songs.

\begin{table}[htbp]
\centering\small
\begin{tabular}{lrrrrr}
\hline
Feature set & 25 & 50 & 100 & 200 & All\\
\hline
S/D/R/P & 82.56 & 90.48 & 75.04 & 84.76 & 76.76\\
S/D/R/P + F & 56.76 & 83.24 & 59.12 & 78.00 & 63.76\\
S/D/R/P + H & 79.16 & 89.56 & 77.20 & 86.56 & 80.08\\
S/D/R/P + M & 83.32 & 90.80 & 78.56 & 88.84 & 83.28\\
All seven & 60.76 & 89.40 & 74.76 & 87.04 & 79.36\\
\hline
\end{tabular}
\caption{Mean Mureka sensitivity (\%) by original training-group cap. Repeated
models share the same external songs; these values are not independent trials.}
\label{tab:mureka-quantity}
\end{table}

Transfer sensitivity is not monotonic in training cap. F lowers the baseline
cap-mean sensitivity at all five caps; M raises it at all five caps. Selecting
a cap using these external results would consume this cohort for development
and require new external confirmation. No such selection is made here.

\subsection{Development gains are not generator-independent guarantees}
On the original source-pair/fold development evaluation, adding F to S/D/R/P
increased mean AUC from 0.906513 to 0.936725 and balanced accuracy from
0.843986 to 0.870106. On Mureka, the corresponding mean sensitivity instead
fell by 13.00 percentage points. The endpoints and populations differ; this
is not an AUC-to-sensitivity numerical delta. Rather, it demonstrates that
the development gain did not imply improved external-generator recall.

These observations are consistent with generator/source dependence but do not
identify its acoustic mechanism. Codec, separation, production style and genre
remain plausible confounds. The conditional M gain merits further testing,
but an AI-only cohort cannot supply Human false-positive rates. Therefore no
standalone Mureka balanced accuracy, ROC AUC, specificity or overall detection
accuracy is reported. Saraga103 preparation is ongoing; its external Human
false-positive evaluation is not yet available. Joint two-class claims require
a separately specified source-weighting protocol. Historical test consumption
and possible performance, excerpt and pretraining overlap remain explicit.

\subsection{Code and artifact references}
The frozen transfer contract is
\path{preregistration/mureka60_frozen_v4_transfer_frozen_v1.json}; parent
pre-score review is \path{audit/mureka60_transfer_draft_parent_review_v1.json}.
Code is \path{code/score_mureka60_frozen_v4.py} and the separately implemented
\path{code/audit_mureka60_transfer_v1.py}. The passed audit is
\path{audit/mureka60_frozen_v4_transfer_numerical_audit_v1.json}, SHA256
\path{47b567d8e126cbcbefa51794b6f2b0af0b462091d719bce1a8fa1fc4ea45da97}.
All per-model metrics and raw predictions are under
\path{results/mureka60_frozen_v4_transfer_results_v1/}; its publication manifest
SHA256 is \path{4c620b612c9ba79042a59799b3aa534bec3a2f993d6dd0e41b6b091181fa7274}.
The complete narrative and predefined overview are
\path{MUREKA60_FROZEN_TRANSFER_RESULTS_EN.md}. Measurement provenance and
missingness are documented in \path{EXTERNAL_MEASUREMENT_CHECKPOINT_EN.md}.
